\documentclass[12pt]{article}

\usepackage[T1]{fontenc}
\usepackage[utf8]{inputenc}
\usepackage{amsmath,amssymb}
\usepackage[version=4]{mhchem}
\usepackage{enumitem}
\usepackage{setspace}


\setstretch{1.0}
\setlist[itemize]{noitemsep, topsep=2pt}
\setlist[enumerate]{noitemsep, topsep=2pt}

\usepackage[
    top=2cm,
    bottom=2cm,
    left=2cm,
    right=2cm
]{geometry}

\usepackage[
    backend=biber,
    style=authoryear,
    doi=true,
    url=false,
    sorting=nyt,
    maxcitenames=2,
    giveninits=true
]{biblatex}

\addbibresource{isotopes.bib}

\setcounter{subsection}{0}
\renewcommand{\thesubsection}{A\arabic{subsection}}

\title{
    \textbf{A closure study to validate a novel isotope-aware
    Lagrangian cloud-microphysics model\\ against airborne measurements}\\[0.4em]
    \large Próba zgodności z pomiarami lotniczymi nowego lagranżowskiego modelu
    mikrofizyki chmur obejmującego zjawiska izotopowe
}
\date{}
\begin{document}

\maketitle

%=========================================================
\section{Scientific goal of the project}
%=========================================================

Cloud microphysical processes such as condensation, evaporation, deposition,
sublimation, collision--coalescence, and freezing are among the dominant sources
of uncertainty in weather and climate prediction \parencite{IPCC_2023}.
These processes
operate at the particle scale and are therefore only indirectly constrained in
current atmospheric models.

Stable water isotopologues (HDO, \ce{H2^{18}O}, \ce{H2^{17}O}) provide a natural
tracer of cloud evolution because isotopic fractionation occurs systematically
during phase transitions and diffusive exchange \parencite{1994_Gedzelman_Arnold}.
As a result, isotopic composition encodes integrated information about the
thermodynamic and microphysical history of cloud systems.

This project investigates whether isotopic observations are sufficient to
constrain cloud microphysics in a particle-resolved framework and whether a
physically consistent \emph{isotopic closure} can be achieved.

\textbf{Research questions:}
\begin{enumerate}
\item Can isotopic closure be achieved in a particle-based cloud model using current observations?
\item Which microphysical processes generate identifiable isotopic signatures?
\item To what extent do isotopic observations constrain supersaturation variability,
phase partitioning, and mixing?
\item What is the minimal physical and numerical complexity required for closure?
\end{enumerate}

We hypothesise that isotopic composition contains sufficient information to constrain
key cloud microphysical pathways, but only when particle-scale variability is resolved.

%=========================================================
\section{Significance of the project}
%=========================================================

Cloud processes remain a leading uncertainty in climate sensitivity and precipitation
prediction \parencite{IPCC_2023}.
Existing parameterisations rely on bulk Eulerian formulations
that cannot resolve stochastic particle-scale variability.

Stable isotopes have been widely used to trace atmospheric moisture transport and
have been incorporated into global models, improving understanding of the hydrological
cycle \parencite{2020_Risi_EtAl, 2016_Galewsky_EtAl}.
However, these models rely on bulk
microphysics and do not explicitly represent particle-scale fractionation processes.

In contrast, particle-based cloud models based on the Super-Droplet Method (SDM)
represent a state-of-the-art framework for cloud microphysics \parencite{2009_Shima_EtAl},
explicitly resolving aerosol, droplet, and ice particle evolution under stochastic
interactions.

Despite these advances, isotopic composition has not been systematically integrated
into particle-resolved cloud modelling.
Consequently, it is still unknown whether
isotopic observations can provide a closure constraint linking observed cloud states
with microphysical process models.

This project introduces \emph{isotopic closure} as a new diagnostic principle in
cloud physics. 
It enables evaluation of cloud models not only against bulk variables
but against full isotopic state variables, providing a fundamentally new constraint
on cloud evolution.

%=========================================================
\section{Concept and work plan}
%=========================================================

The project is structured into three coupled components:

%-------------------------
\subsection{Theory: isotopic cloud microphysics}
%-------------------------

A unified theoretical framework will be developed for isotope fractionation in
vapour--liquid--ice systems, including equilibrium and kinetic processes.
A key element is the Bolin number formulation
\parencite{1958_Bolin}, extended here to mixed-phase cloud systems as a control parameter
for isotopic equilibration.
The framework identifies limiting regimes such as equilibrium fractionation,
Rayleigh distillation, and diffusion-limited growth, and provides a closed
mathematical formulation of isotopic evolution in clouds.
Additionally, uncertainty propagation and sensitivity analysis of isotopic evolution
with respect to microphysical parameters are included.

%-------------------------
\subsection{Numerics: particle-based implementation}
%-------------------------

The theoretical framework will be implemented in the Super-Droplet Method (SDM)
(Shima et al., 2009) using the PySDM environment (Bartman et al., 2022; De Jong et al., 2023).
Each computational particle will carry prognostic isotopic composition in vapour,
liquid, and ice phases.
The model explicitly represents isotope fractionation during condensation,
evaporation, deposition, sublimation, and freezing, fully coupled to particle
dynamics.
Simulations will be performed in idealised configurations (box, parcel, and 1D column)
to isolate governing processes and quantify sensitivity to supersaturation variability,
droplet size distribution, and phase partitioning.

%-------------------------
\subsection{Data: observational and experimental validation}
%-------------------------

Model predictions will be evaluated using airborne isotopic measurements
(e.g. CAESAR campaign) and laboratory experiments in collaboration with NCAR.
These datasets provide complementary constraints across vapour, cloud water,
ice, and precipitation phases.
The analysis focuses on:
(i) agreement between simulated and observed isotopic composition,
(ii) identification of microphysical processes controlling isotopic signals,
and (iii) assessment of whether isotopic observations enable closure.

%-------------------------
\subsection*{Integration and preliminary results}
%-------------------------

Preliminary work demonstrates feasibility of isotope tracking in SDM and introduces
a dimensionless Bolin number describing isotopic equilibration regimes. The
implementation has been validated against analytical fractionation limits
\parencite{1994_Gedzelman_Arnold}.

%-------------------------
\subsection*{Milestones}
%-------------------------

\begin{itemize}
\item \textbf{M1:} Theoretical isotope cloud microphysics framework
\item \textbf{M2:} Isotope-enabled SDM implementation and validation
\item \textbf{M3:} Integration with observational and laboratory datasets
\item \textbf{M4:} Isotopic closure assessment and minimal model identification
\end{itemize}

%=========================================================
\section{Research methodology}
%=========================================================

The project combines theoretical analysis, particle-based numerical simulation,
and systematic data--model comparison.
The core framework is the Super-Droplet Method (SDM), which represents clouds
as ensembles of computational particles resolving stochastic cloud microphysics.
The model will be extended to include prognostic isotopic
composition for all phases.
This enables explicit simulation of isotope fractionation at the particle level,
fully coupled to condensation, evaporation, collision--coalescence, and ice-phase
processes.
Simulations in idealised configurations will isolate key mechanisms controlling
isotopic evolution and quantify sensitivity to environmental forcing.
Model output will be analysed using statistical diagnostics, sensitivity analysis,
and uncertainty quantification, including ensemble simulations.
Observational constraints from aircraft campaigns (CAESAR) and laboratory experiments
will be used for model evaluation and isotopic closure assessment.
The workflow is implemented in Python with reproducible Jupyter notebooks and
continuous integration to ensure numerical consistency.

%=========================================================
\section{Composition and qualifications of the research team}
%=========================================================

\noindent\textbf{Principal Investigator:} PhD student in physical sciences with background
in mathematical modelling and scientific computing.
Developer of isotope-enabled
SDM framework and experienced in particle-based cloud modelling.

\noindent\textbf{Scientific mentor:} Expert in Lagrangian cloud microphysics and co-developer
of PySDM (Shima et al., 2009; Bartman et al., 2022; De Jong et al., 2023).

\noindent\textbf{Student (MSc/Erasmus+/PhD):} Supports development of simplified isotope-enabled
models and contributes to simulation and analysis workflows.

The team is embedded in a research group specialising in atmospheric isotope systematics
and tracer-based hydrological processes.

%=========================================================
\appendix

\clearpage
\printbibliography

\end{document}